% Plantilla creada por Sebastian Manrique.

% La siguiente plantilla fue creada con el fin de que se pueda usar para
% la presentación de trabajos y que puedas tener una mejor carta de
% presentación frente a tus evaluadores. :D

% Puedes usarla con los fines que desees sin embargo no olvides mencionar
% de donde la obtuviste cuando la socialices con tus colegas. 

% Para conocer todas las herramientas que poseen puedes leer el archivo
% README.md que probablemente esté en la parte inferior o si clonaste el
% repositorio, lo encuentras en la misma carpeta en la cual está este archivo.

\documentclass[letterpaper,11pt]{article}

\usepackage[spanish,es-tabla]{babel}
\usepackage[top=2cm,bottom=2cm,right=2cm,left=2cm]{geometry}
\usepackage{amsmath}
\usepackage{amssymb}
\usepackage{amsthm}
\usepackage{empheq}
\usepackage{mdframed}
\usepackage{booktabs}
\usepackage{lipsum}
\usepackage{graphicx}
\usepackage{color}
\usepackage{psfrag}
\usepackage{pgfplots}
\usepackage{bm}
\usepackage{multicol}
\usepackage{wrapfig}
\usepackage{float}
\usepackage{wrapfig}
\usepackage{multirow}
\usepackage{biblatex}
\usepackage{csquotes}
\usepackage{parskip}
\usepackage{cancel}
\usepackage{listing}
\usepackage{tikz}
\usepackage{verbatim}
\usepackage{subcaption}
\usepackage{enumitem}
\usepackage{xcolor}
\usepackage{xfrac}
\usepackage[load-configurations = abbreviations]{siunitx}
\usepackage{fix-cm}
\usepackage{yhmath}
\usepackage[bookmarks=true,colorlinks=true,linkcolor=black,citecolor=black,menucolor=black,urlcolor=black,hidelinks,breaklinks=true]{hyperref}
\usepackage{mathrsfs}
\usepackage{tcolorbox}
\usepackage{amsfonts}
\usepackage[bb=ams]{mathalpha}
\usepackage{mathrsfs}
\usepackage{xstring}

\documentclass[letterpaper,11pt]{article}

\usepackage[spanish]{babel}
\usepackage[top=2cm,bottom=2cm,right=2cm,left=2cm]{geometry}
\usepackage{amssymb}
\usepackage{amsmath}
\usepackage{amsthm}
\usepackage{empheq}
\usepackage{mdframed}
\usepackage{booktabs}
\usepackage{lipsum}
\usepackage{graphicx}
\usepackage{color}
\usepackage{psfrag}
\usepackage{pgfplots}
\usepackage{bm}
\usepackage{multicol}
\usepackage{wrapfig}
\usepackage{float}
\usepackage{wrapfig}
\usepackage{multirow}
\usepackage{biblatex}
\usepackage{csquotes}
\usepackage{parskip}
\usepackage{cancel}
\usepackage{listing}
\usepackage{tikz}
\usepackage{verbatim}
\usepackage{subcaption}
\usepackage{enumitem}
\usepackage{xcolor}
\usepackage{xfrac}

\graphicspath{{fig/}}
\addbibresource{biblio.bib}
\usepgfplotslibrary{colorbrewer}
\pgfplotsset{width=8cm,compat=1.9}
\definecolor{negro}{RGB}{0,0,0}
\definecolor{azul}{RGB}{0,0,255}
\definecolor{verde}{RGB}{0,255,0}
\definecolor{cian}{RGB}{0,255,255}
\definecolor{rojo}{RGB}{255,0,0}
\definecolor{magenta}{RGB}{255,0,255}
\definecolor{amarillo}{RGB}{255,255,0}
\definecolor{blanco}{RGB}{255,255,255}

\definecolor{limegreen}{cmyk}{0.75,0,1,0}
\definecolor{orange}{cmyk}{0,0.8,0.95,0}
\definecolor{purple}{cmyk}{0.75,1,0,0}

\begin{document}

\onecolumn\begin{@twocolumntrue}
    \begin{minipage}{0.3\textwidth}{\includegraphics[scale=0.24]{Escudo_UD.pdf}} % Cambiar escudo
    \end{minipage}
    \vspace{10pt}
    \begin{minipage}{0.677\textwidth}
        \begin{center}
            \vspace{12mm}
        
            \Large{\textbf{Tarea 3: osciladores y ecuaciones diferenciales}}
            \vspace{3mm}
            
            \large{\textbf{Juan Sebastian Manrique Moreno}}
            \vspace{2mm}
    
            \large{\textit{Ecuaciones diferenciales, Programa Académico de Física, Universidad Distrital Francisco José de Caldas}}
            \vspace{1mm}
            
            Octubre de 2022 % Formato (mes) de (año)
        \end{center}
\vspace{5pt}
\end{minipage}

\centerline{\rule{0.95\textwidth}{0.4pt}} % Línea negra
\end{@twocolumntrue}

%\begin{center}
    \Large \textbf{Resumen}
    
    \small
    <<Texto>>
\end{center}
\hspace{10pt}\textbf{Palabras clave:} <<Texto>>
\begin{center}
    \Large \textbf{Abstract}
    
    \small
    <<Text>>
\end{center}
\hspace{10pt}\textbf{Keywords:} <<Text>>

\centerline{\rule{0.95\textwidth}{0.4pt}}
\vspace{15pt}

\begin{resumen}
    La maravilla que es esto
\end{resumen}

\begin{abs}
    Buenas tardes
\end{abs}

\begin{multicols}{2}

\section{Introducción}

\section{Sección 1} 

\begin{equation}
    Ecuaciones
\label{ecuacion}
\end{equation}

\begin{equation*} % Para ecuaciones sin número y que el igual se mantenga en la misma sangría
\begin{split}
    texto1 & = texto---2 \\
    texto---2 & = texto------3 \\
    texto------3 & = texto4
\end{split}
\end{equation*}

\section{Sección 2}

<<Texto>>

\end{multicols}

\subsection{Subsección 1}

<<Variación entre dos y una columna>> -------------------------------------------------------------------

\begin{multicols}{2}

\subsection{Subsección 2}

<<Texto>>

\subsubsection{Subsubsección 1}.

<<Texto>>

\begin{figure}[H]
    \centering
    \includegraphics[scale=0.32]{pollito.png}
    \caption{Inserción de una figura}
    \label{figura}
\end{figure}

\section{Conclusiones}

<<Conclusión>>

\end{multicols}

\section{Anexos}

<<Anexo>>

\end{document}